\documentclass[11pt,oneside]{report}
\usepackage[margin=0.8in]{geometry}
\usepackage{fontspec}
\usepackage{graphicx}
\usepackage{xcolor}
\usepackage{booktabs}
\usepackage{longtable}
\usepackage{array}
\usepackage{enumitem}
\usepackage{hyperref}
\usepackage{listings}
\usepackage{caption}
\usepackage{float}
\usepackage{titlesec}
\usepackage{fancyhdr}
\usepackage{lastpage}

\setmainfont{TeX Gyre Pagella}
\setsansfont{TeX Gyre Heros}
\setmonofont{Menlo}[Scale=0.82]

\definecolor{navy}{HTML}{172033}
\definecolor{steel}{HTML}{344054}
\definecolor{blue}{HTML}{1F6FEB}
\definecolor{lightblue}{HTML}{EDF4FF}
\definecolor{lightgray}{HTML}{F2F4F7}
\definecolor{codebg}{HTML}{F7F9FC}

\hypersetup{
  colorlinks=true,
  linkcolor=blue,
  urlcolor=blue,
  citecolor=blue,
  pdftitle={Precision Dot Product DV Lab: Final Verification Report},
  pdfauthor={Sasha Katne}
}

\graphicspath{{doc/final_report/assets/figures/}{doc/final_report/assets/waveforms/}}
\setlist[itemize]{leftmargin=*, itemsep=2pt, topsep=3pt}
\setlist[enumerate]{leftmargin=*, itemsep=2pt, topsep=3pt}
\renewcommand{\arraystretch}{1.18}
\titleformat{\chapter}{\huge\bfseries\sffamily\color{navy}}{\thechapter}{0.75em}{}
\titleformat{\section}{\Large\bfseries\sffamily\color{navy}}{\thesection}{0.75em}{}
\titleformat{\subsection}{\large\bfseries\sffamily\color{steel}}{\thesubsection}{0.75em}{}

\lstdefinestyle{sv}{
  language=Verilog,
  basicstyle=\ttfamily\footnotesize,
  backgroundcolor=\color{codebg},
  frame=single,
  rulecolor=\color{lightgray},
  breaklines=true,
  columns=fullflexible,
  keepspaces=true,
  keywordstyle=\color{blue}\bfseries,
  commentstyle=\color{steel},
  showstringspaces=false
}
\lstdefinestyle{shell}{
  basicstyle=\ttfamily\footnotesize,
  backgroundcolor=\color{codebg},
  frame=single,
  rulecolor=\color{lightgray},
  breaklines=true,
  columns=fullflexible,
  keepspaces=true,
  showstringspaces=false
}

\pagestyle{fancy}
\setlength{\headheight}{14pt}
\fancyhf{}
\lhead{\sffamily Precision Dot Product DV Lab}
\rhead{\sffamily Final Verification Report}
\cfoot{\sffamily \thepage\ of \pageref{LastPage}}

\newcommand{\statuspass}{\textcolor{blue}{\textbf{PASS}}}
\newcommand{\statusfail}{\textcolor{red}{\textbf{FALSIFIED}}}
\newcommand{\artifactnote}{Raw tool logs, waveform databases, UCDB files, and transcripts are generated artifacts. They are not checked into the repository; this report includes summarized evidence and curated figures only.}

\begin{document}

\begin{titlepage}
  \centering
  \vspace*{0.6in}
  {\Huge\bfseries\sffamily\color{navy} Precision Dot Product DV Lab\par}
  \vspace{0.16in}
  {\Large\bfseries Final Verification Report\par}
  \vspace{0.28in}
  {\large INT8, BF16, and NVFP4 Dot-Product Formal + UVM Verification\par}
  \vspace{0.45in}
  \begin{tabular}{rl}
    \textbf{Author:} & Sasha Katne \\
    \textbf{Project:} & nvfp4-dotprod-formal-dv \\
    \textbf{Report date:} & 2026-07-03 \\
    \textbf{Status:} & Complete sign-off, public-safe report \\
  \end{tabular}
  \vfill
  \begin{minipage}{0.82\textwidth}
  \centering
  This document consolidates the problem statement, DUT architecture,
  verification strategy, formal properties, regenerated evidence, waveform
  figures, coverage closure, and conclusions for the mixed-precision
  dot-product verification project.
  \end{minipage}
  \vfill
\end{titlepage}

\tableofcontents
\listoffigures
\listoftables
\clearpage

\chapter*{Executive Summary}
\addcontentsline{toc}{chapter}{Executive Summary}

This project verifies an 8-lane mixed-precision vector dot-product unit across
three arithmetic tiers: INT8, BF16, and NVIDIA-style NVFP4 block-scaled FP4.
The verification strategy combines direct formal equivalence, assume-guarantee
formal decomposition, directed simulation, UVM constrained-random testing, code
and functional coverage, and bug-injection checks.

\begin{table}[H]
\centering
\caption{Final sign-off snapshot}
\begin{tabular}{ll}
\toprule
\textbf{Area} & \textbf{Final status} \\
\midrule
Formal proof jobs & 22 total: 11 clean and 11 bug-injected \\
Clean formal jobs & All target assertions proven (BF16 proven over the full input space) \\
Bug-injected formal jobs & Every mutation falsified at least one intended property \\
UVM tests & 7 tests, zero mismatches, zero leftovers, zero fatal/errors \\
Coverage & 100.00\% reachable merged coverage \\
Waivers & One structurally unreachable branch leg, documented and justified \\
\bottomrule
\end{tabular}
\end{table}

The central verification principle is a single-source-of-truth golden model:
SystemVerilog reference functions in \texttt{ref/} are used by both formal
properties and simulation scoreboards. This avoids the common verification
failure mode where a formal checker and a simulation checker silently validate
against different reference behavior.

\artifactnote

\chapter{Problem Statement and Project Introduction}

\section{Problem Statement}

Modern accelerator datapaths increasingly mix integer and compressed floating
formats in the same compute block. A dot-product unit is deceptively small, but
it is verification-dense: sign handling, exact accumulation, final rounding,
special values, packed element formats, block scaling, ready/valid protocol, and
coverage closure all interact.

The problem addressed by this project is:

\begin{quote}
Prove and test that a unified vector dot-product design computes correct results
for INT8, BF16, and NVFP4 modes, while preserving a correct streaming
ready/valid interface and demonstrating that the verification environment catches
representative injected bugs.
\end{quote}

\section{Verification Goals}

\begin{itemize}
  \item Prove arithmetic equivalence for INT8 exhaustively.
  \item Verify the sequential wrapper's ready/valid and backpressure behavior.
  \item Prove BF16 behavior using a decomposition that avoids monolithic
        nonlinear equivalence bottlenecks.
  \item Prove NVFP4 behavior using pure-DUT-net top assertions and standalone
        lane, scale, and final-round proofs.
  \item Use the same reference functions from formal and UVM scoreboards.
  \item Demonstrate proof strength with bug-injected falsification runs.
  \item Close reachable coverage to 100.00\% and document any unreachable logic.
\end{itemize}

\section{Numeric Modes}

\begin{longtable}{p{0.16\textwidth}p{0.33\textwidth}p{0.39\textwidth}}
\caption{Numeric tiers verified in the project}\\
\toprule
\textbf{Mode} & \textbf{Computation} & \textbf{Verification emphasis} \\
\midrule
\endfirsthead
\toprule
\textbf{Mode} & \textbf{Computation} & \textbf{Verification emphasis} \\
\midrule
\endhead
INT8 & Eight signed 8-bit products summed exactly and cast to a 32-bit output. &
Exhaustive top equivalence; saturation shown structurally unreachable for one
8-lane dot product. \\
BF16 & Eight BF16 products accumulated exactly in the fixed-point window
[119,134], then rounded once to FP32; out-of-window operands are flagged as an
invalid operation. & Lane decode/product, alignment, special case ladder, and
final round proven in separate formal pieces; equivalence proven over the full
BF16 input space. \\
NVFP4 & Sixteen E2M1 elements per vector share UE4M3 block scales; result is
FP32. & Lane product, scale multiply, inner sum, and final normalization
verified with assume-guarantee composition. \\
\bottomrule
\end{longtable}

\chapter{DUT Architecture}

\section{Unified Top-Level Datapath}

The top-level DUT exposes widened 16-bit operand lanes and selects the active
numeric tier using \texttt{mode}. INT8 consumes the low byte; BF16 consumes the
full 16-bit lane; NVFP4 interprets the operand lanes as a packed block of sixteen
4-bit elements plus UE4M3 scale bytes.

\begin{figure}[H]
\centering
\includegraphics[width=0.94\textwidth]{architecture_dut.pdf}
\caption{Unified DUT datapath}
\end{figure}

The output port is 32 bits. In INT8 mode it carries the sign-extended integer
result. In BF16 and NVFP4 modes it carries an IEEE binary32 bit pattern. Status
fields carry saturation, invalid, NaN, and infinity indicators; the legacy
\texttt{sat} output aliases \texttt{status.sat}.

\section{Sequential Wrapper}

The sequential wrapper adds a two-stage ready/valid pipeline around the
combinational core. Each beat carries a valid bit that shifts through the two
stages in lockstep with its data (\texttt{in\_valid \& in\_ready} into
\texttt{v\_s1}, then into \texttt{out\_valid}), so this is a pipeline rather
than a control FSM: there is no state register and no next-state logic, only a
single valid bit advancing down the pipe. A global stall, derived from a single
\texttt{always\_ff} enable, freezes both stages together whenever a valid output
is present and the consumer deasserts \texttt{out\_ready}, so a held beat is
never dropped, reordered, or partially updated. The key protocol contract is
that a held output must remain valid and stable under backpressure.

\begin{figure}[H]
\centering
\includegraphics[width=0.94\textwidth]{architecture_seq_dut.pdf}
\caption{Sequential ready/valid wrapper: a two-stage pipeline with a valid bit
shifting in lockstep with its data, and one stall enable freezing both stages
(not a control FSM).}
\end{figure}

\section{BF16 Datapath}

The BF16 datapath decodes operands, flushes subnormal inputs to zero, multiplies
lanes, aligns products into a fixed-point accumulator window, resolves IEEE
special cases, and performs one final round-nearest ties-to-even conversion to
FP32.

\begin{figure}[H]
\centering
\includegraphics[width=0.94\textwidth]{architecture_bf16_dut.pdf}
\caption{BF16 datapath and proof decomposition}
\end{figure}

\section{NVFP4 Datapath}

NVFP4 mode factors the block-scaled dot product into an exact inner sum and a
scale product:
\[
  \sum_i (a_i s_A)(b_i s_B) = (s_A s_B) \sum_i (a_i b_i)
\]
The inner E2M1 product sum is exact, the UE4M3 scale product is checked
separately, and final normalization produces FP32.

\begin{figure}[H]
\centering
\includegraphics[width=0.94\textwidth]{architecture_nvfp4_dut.pdf}
\caption{NVFP4 block-scaled datapath}
\end{figure}

\chapter{Verification Plan}

\section{Strategy Overview}

The verification plan intentionally uses different proof shapes for different
logic classes. Linear INT8 arithmetic is handled as a direct exhaustive
equivalence proof. Nonlinear BF16 and NVFP4 logic is split into smaller proofs
that can converge independently. The sequential wrapper is verified with
protocol SVA and UVM stress tests.

\begin{figure}[H]
\centering
\includegraphics[width=0.94\textwidth]{architecture_uvm_tb.pdf}
\caption{UVM environment architecture}
\end{figure}

\begin{table}[H]
\centering
\caption{Verification methods by target}
\begin{tabular}{p{0.18\textwidth}p{0.38\textwidth}p{0.32\textwidth}}
\toprule
\textbf{Target} & \textbf{Formal plan} & \textbf{Simulation plan} \\
\midrule
INT8 top & Direct top equivalence against \texttt{dotprod\_ref}. &
Directed corner tests plus UVM random/corner tests. \\
Sequential wrapper & Ready/valid protocol assertions. &
Random backpressure and deterministic waveform case. \\
BF16 & Lane, align, round, special, and top assume-guarantee proofs. &
BF16 random and directed special-value tests. \\
NVFP4 & Lane, scale, final round, and pure-DUT-net top proof. &
NVFP4 random and directed block-scale tests. \\
\bottomrule
\end{tabular}
\end{table}

\section{UVM Scoreboard And Coverage}

The UVM environment uses input and output agents around the ready/valid wrapper.
The scoreboard queues accepted input transactions, calls the appropriate golden
reference based on \texttt{mode}, and compares each accepted output in order.
Coverage tracks value classes, protocol states, BF16 operand/result classes, and
NVFP4 element/scale/result classes.

\section{Bug-Injection Plan}

Bug-injection runs are not optional smoke tests; they are the proof-suite
sensitivity check. Each mutation is expected to falsify a specific property. A
clean proof without a corresponding falsification can still be vacuous or too
weak, so this project pairs clean proof runs with targeted failing variants.

\chapter{Formal Properties}

\section{INT8 Top Equivalence}

The INT8 proof constrains the top to INT8 mode, evaluates the low byte of each
operand lane with \texttt{dotprod\_ref}, and asserts that result, status, and the
\texttt{sat} alias all match the reference.

\begin{lstlisting}[style=sv,caption={INT8 top equivalence property}]
am_int8_mode: assume (mode == FMT_INT8);

a_result_matches_ref: assert (
    result == FP32_W'(ref_r) &&
    status.sat == ref_sat &&
    sat == ref_sat &&
    !status.invalid && !status.is_nan && !status.is_inf);

c_int8_sat_unreachable: cover (status.sat == 1'b1);
c_positive_result:      cover ($signed(result) > 0);
c_negative_result:      cover ($signed(result) < 0);
\end{lstlisting}

\section{Sequential Protocol Properties}

\begin{lstlisting}[style=sv,caption={Sequential ready/valid protocol properties}]
p_hold_stable: assert property
  (out_valid && !out_ready |=>
     out_valid && $stable(result) && $stable(status) && $stable(sat));

p_no_out_at_reset: assert property
  ($rose(rst_n) |-> !out_valid);

p_stall_blocks_ready: assert property
  (out_valid && !out_ready |-> !in_ready);

p_sat_alias: assert property (sat == status.sat);

c_backpressure: cover property (out_valid && !out_ready);
c_accept:       cover property (in_valid && in_ready);
\end{lstlisting}

\section{BF16 Properties}

BF16 is decomposed into lane, align, round, special ladder, and top
assume-guarantee checks. The top proof blackboxes the lane multiplier, assumes
the standalone lane guarantee, and proves linear pre-round reduction plus
special case equivalence.

\begin{lstlisting}[style=sv,caption={BF16 top assume-guarantee properties}]
am_bf16_mode: assume (mode == FMT_BF16);

a_acc_is_lane_sum:            assert (sum_bf16 == lane_sum);
a_special_valid_matches_ref:  assert (special_valid  == ref_pr.special_valid);
a_special_result_matches_ref: assert (special_result == ref_pr.special_result);
a_special_status_matches_ref: assert (special_status == ref_pr.special_status);

c_numeric:  cover (!ref_pr.special_valid);
c_special:  cover (ref_pr.special_valid);
c_nan:      cover (ref_pr.special_status.is_nan);
c_invalid:  cover (ref_pr.special_status.invalid);
\end{lstlisting}

\section{NVFP4 Properties}

The NVFP4 top proof avoids placing a 16-lane golden multiply inside the hardest
miter. Instead, it assumes the independently proven lane guarantee and asserts
only over DUT nets for alignment and accumulation.

\begin{lstlisting}[style=sv,caption={NVFP4 top assume-guarantee properties}]
am_lane_ref:   assume (nvfp4_prod[gi] == ref_mul_nvfp4(nv_ea[gi], nv_eb[gi]));
a_align_lane:  assert (nvfp4_wide[gi] == NVFP4_INNER_W'(nvfp4_prod[gi].prod));

am_nvfp4_mode: assume (mode == FMT_NVFP4);
a_inner_is_lane_sum: assert (nvfp4_inner_sum == lane_sum);
a_scale_sig_ref:     assert (nvfp4_scale_sig == ref_pr.scale_sig);
a_scale_exp_ref:     assert (nvfp4_scale_exp == ref_pr.scale_exp);
a_scale_nan_ref:     assert (nvfp4_scale_is_nan == ref_pr.scale_is_nan);
\end{lstlisting}

The final-round proof checks both normal numeric behavior and NaN-scale bypass:

\begin{lstlisting}[style=sv,caption={NVFP4 final-round properties}]
a_numeric_round: assert (scale_is_nan ||
  (result == ref_result && status == ref_status));

a_nan_bypass: assert (!scale_is_nan ||
  (result == FP32_QNAN && status.invalid && status.is_nan &&
   !status.is_inf && !status.sat));
\end{lstlisting}

\chapter{Sign-Off Evidence and Regenerated Waveforms}

\section{Formal Results}

The formal, UVM, and coverage tables consolidate the M5 sign-off record
preserved in the repository documentation. The timing figures in this chapter
were regenerated with report-specific self-checking testbenches so the PDF
contains curated waveform evidence without checking raw VCD or transcript
artifacts into source control.

\begin{longtable}{p{0.36\textwidth}p{0.16\textwidth}p{0.18\textwidth}p{0.18\textwidth}}
\caption{Unified formal proof results}\\
\toprule
\textbf{Proof job} & \textbf{Tier} & \textbf{Clean outcome} & \textbf{Bug outcome} \\
\midrule
\endfirsthead
\toprule
\textbf{Proof job} & \textbf{Tier} & \textbf{Clean outcome} & \textbf{Bug outcome} \\
\midrule
\endhead
\texttt{fpv\_run\_top} & INT8 & 1 assert proven & Bug falsifies result equivalence \\
\texttt{fpv\_run\_seq} & INT8 wrapper & 4 asserts proven & Bug falsifies \texttt{p\_hold\_stable} \\
\texttt{fpv\_run\_lane\_bf16} & BF16 & 24 asserts proven & FTZ and out-of-range mutations falsify \\
\texttt{fpv\_run\_align\_bf16} & BF16 & 8 asserts proven & Align mutation falsifies \\
\texttt{fpv\_run\_round\_bf16} & BF16 & 2 asserts proven & Round mutation falsifies \\
\texttt{fpv\_run\_special\_bf16} & BF16 & 5 asserts proven & Covered by top special mutation \\
\texttt{fpv\_run\_bf16\_top} & BF16 & 4 asserts proven & Special mutation falsifies \\
\texttt{fpv\_run\_lane\_nvfp4} & NVFP4 & 5 asserts proven & Lane mutation falsifies \\
\texttt{fpv\_run\_scale\_nvfp4} & NVFP4 & 3 asserts proven & Scale mutation falsifies \\
\texttt{fpv\_run\_round\_nvfp4} & NVFP4 & 2 asserts proven & NaN-bypass mutation falsifies \\
\texttt{fpv\_run\_nvfp4\_top} & NVFP4 & 20 asserts proven & NaN-detect mutation falsifies \\
\bottomrule
\end{longtable}

\section{UVM And Coverage Results}

\begin{table}[H]
\centering
\caption{Final UVM regression tests}
\begin{tabular}{lrrrr}
\toprule
\textbf{Test} & \textbf{Items} & \textbf{Matched} & \textbf{Mismatched} & \textbf{Leftover} \\
\midrule
\texttt{dotprod\_random\_test} & 500 & 500 & 0 & 0 \\
\texttt{dotprod\_backpressure\_test} & 1000 & 1000 & 0 & 0 \\
\texttt{dotprod\_corner\_test} & 9 & 9 & 0 & 0 \\
\texttt{dotprod\_bf16\_test} & 500 & 500 & 0 & 0 \\
\texttt{dotprod\_bf16\_corner\_test} & 10 & 10 & 0 & 0 \\
\texttt{dotprod\_nvfp4\_test} & 500 & 500 & 0 & 0 \\
\texttt{dotprod\_nvfp4\_corner\_test} & 8 & 8 & 0 & 0 \\
\bottomrule
\end{tabular}
\end{table}

\begin{table}[H]
\centering
\caption{Merged coverage closure}
\begin{tabular}{lr}
\toprule
\textbf{Metric} & \textbf{Result} \\
\midrule
Statements & 100\% \\
Branches & 100\% \\
Expressions & 100\% \\
Covergroups & 100\% \\
Covergroup bins & 206/206 \\
Total & 100.00\% \\
\bottomrule
\end{tabular}
\end{table}

\section{Waveform: Combinational Top Behavior}

Figure \ref{fig:top-modes-wave} shows the report-specific deterministic top
waveform. The sequence exercises INT8, BF16, NVFP4 numeric, and NVFP4 NaN-scale
cases. The DUT result tracks the expected result from the same golden reference
functions used by the scoreboards.

\begin{figure}[H]
\centering
\includegraphics[width=\textwidth]{top_modes.pdf}
\caption{Combinational top waveform: multi-mode behavior}
\label{fig:top-modes-wave}
\end{figure}

\section{Waveform: Clean Backpressure Hold}

Figure \ref{fig:seq-clean-wave} shows the sequential wrapper holding output
valid and result stable while \texttt{out\_ready} is low. The report assertion
flag remains low.

\begin{figure}[H]
\centering
\includegraphics[width=\textwidth]{seq_clean.pdf}
\caption{Sequential wrapper waveform: clean backpressure hold}
\label{fig:seq-clean-wave}
\end{figure}

\section{Waveform: Bug-Injected Assertion Failure}

Figure \ref{fig:seq-bug-wave} shows the bug-injected wrapper dropping
\texttt{out\_valid} under backpressure. This directly violates the
\texttt{p\_hold\_stable} contract and raises the report assertion failure flag.

\begin{figure}[H]
\centering
\includegraphics[width=\textwidth]{seq_bug.pdf}
\caption{Sequential wrapper waveform: bug-injected \texttt{p\_hold\_stable} failure}
\label{fig:seq-bug-wave}
\end{figure}

\chapter{Bug-Injection Traceability}

\begin{longtable}{p{0.14\textwidth}p{0.32\textwidth}p{0.28\textwidth}p{0.14\textwidth}}
\caption{Seeded bugs and detection points}\\
\toprule
\textbf{Tier} & \textbf{Seeded fault} & \textbf{Catching property/check} & \textbf{Outcome} \\
\midrule
\endfirsthead
\toprule
\textbf{Tier} & \textbf{Seeded fault} & \textbf{Catching property/check} & \textbf{Outcome} \\
\midrule
\endhead
INT8 & LSB corruption in final rounding & \texttt{a\_result\_matches\_ref} & Falsified \\
INT8 wrapper & Drop held output under stall & \texttt{p\_hold\_stable} & Falsified \\
BF16 & Subnormal no longer flushed to zero & Lane decode/product assertions & Falsified \\
BF16 & Out-of-window operand not flagged & \texttt{a\_product\_matches\_ref}, \texttt{a\_oor\_forces\_invalid\_nan} & Falsified \\
BF16 & Alignment off by one & \texttt{a\_align\_lane} & Falsified \\
BF16 & Truncate instead of RNE & \texttt{a\_numeric\_round} & Falsified \\
BF16 & Inf-minus-Inf dropped & Top special-ladder assertions & Falsified \\
NVFP4 & E2M1 6.0 decoded incorrectly & Lane decode/product assertions & Falsified \\
NVFP4 & UE4M3 exponent bias off by one & \texttt{a\_exp} & Falsified \\
NVFP4 & NaN scale detection threshold changed & \texttt{a\_scale\_nan\_ref} & Falsified \\
NVFP4 & NaN bypass disabled in final round & \texttt{a\_nan\_bypass} & Falsified \\
\bottomrule
\end{longtable}

\chapter{Conclusion}

The project demonstrates a complete verification story for a compact but
high-risk mixed-precision arithmetic block. The final evidence supports these
conclusions:

\begin{itemize}
  \item The INT8 datapath is exhaustively equivalent to the golden reference.
  \item The sequential wrapper satisfies ready/valid stability, reset, stall,
        and status-alias properties.
  \item The BF16 tier is proven by decomposing nonlinear logic into tractable
        standalone proofs and a top-level pre-round assume-guarantee proof, over
        the full BF16 input space: out-of-window operands are detected and
        resolved to a canonical invalid QNaN rather than a silently-wrong result.
  \item The NVFP4 tier is proven using lane, scale, round, and pure-DUT-net top
        properties that avoid expensive golden multiply elaboration in the top
        miter.
  \item The UVM regression validates integrated behavior across all supported
        modes with zero scoreboard residue or mismatches.
  \item Bug-injection runs confirm the proof suite has teeth: every targeted
        mutation is caught.
  \item Reachable coverage closes to 100.00\%; the remaining waiver is
        structurally justified.
\end{itemize}

The resulting methodology is reusable: keep golden references shared, isolate
nonlinear blocks, prove protocol separately from arithmetic, and pair every
clean proof family with a meaningful failing mutation.

\appendix
\chapter{Reproduction Commands}

The report artifacts are generated with the following commands from the repo
root, assuming compatible EDA tools are available in the shell environment:

\begin{lstlisting}[style=shell,caption={Report evidence and PDF build commands}]
bash doc/final_report/scripts/run_report_regression.sh
vsim -c -do doc/final_report/scripts/run_report_waveforms.do
bash doc/final_report/scripts/build_assets.sh
bash doc/final_report/scripts/build_report.sh
\end{lstlisting}

\artifactnote

\end{document}
